\begin{textsample}{POS Dim 2 – persona\_gpt – Score 57.00 – t185\_gpt.txt}  \label{ex:f2_pos_018}
\textbf{Women} and \textbf{gender} non-conforming \textbf{people} are on the \textbf{frontlines} of the \textbf{climate} \textbf{crisis}, yet are too often \textbf{pushed} to the margins of \textbf{environmental} decision-making.
Their work is crucial : defending \textbf{forests} and oceans, \textbf{challenging} fossil \textbf{fuels}, confronting plastic pollution, and insisting that \textbf{climate} \textbf{justice} must include \textbf{justice} for disabled, Indigenous, and trans \textbf{communities}. \textbf{Centering} their \textbf{leadership} is not a matter of representation alone—it is essential to building effective and fair \textbf{climate} solutions.
A key \textbf{barrier} to this is eco-ableism : \textbf{environmental} \textbf{narratives} and \textbf{policies} that ignore, exclude, or even harm disabled \textbf{people}.
When \textbf{climate} and \textbf{environmental} work is designed without disabled \textbf{communities} in mind—whether through inaccessible public actions, one-size-fits-all “ green ” \textbf{policies}, or ignoring the needs of those most affected by extreme weather—it creates deep distrust.
Disabled \textbf{communities} are not just vulnerable to \textbf{climate} \textbf{impacts} ; they are also experts in navigating \textbf{systems} that often fail them.
Overcoming eco-ableism \textbf{means} actively shifting practices and \textbf{priorities} so that disabled \textbf{people} are fully included in \textbf{climate} action, from planning to \textbf{leadership}.
Around the \textbf{world}, \textbf{women} and \textbf{gender} non-conforming activists are already showing what inclusive, community-rooted \textbf{climate} action looks like.
Archana Soreng works to protect \textbf{forests} and uplift Indigenous \textbf{knowledge}, demonstrating how traditional practices and lived experience can guide more sustainable and just \textbf{environmental} \textbf{policies}.
Her work highlights that Indigenous \textbf{peoples} are not simply “ stakeholders ” but rights-holders and \textbf{guardians} of \textbf{ecosystems}, and that their \textbf{knowledge} is \textbf{central} to \textbf{climate} solutions.
In South Korea, Eunbin Kang is part of a \textbf{youth} movement resisting \textbf{coal}.
By confronting \textbf{coal} \textbf{power} and the political \textbf{forces} that sustain it, Eunbin and fellow young activists are \textbf{challenging} a fossil \textbf{fuel} \textbf{system} that drives both \textbf{climate} \textbf{breakdown} and local pollution, and insisting that young \textbf{people} have a say in the energy \textbf{future} they will inherit.
In the Philippines, Hasminah D.
Paudac-Tawano is using the law to hold polluters to account.
Legal battles are a powerful \textbf{tool} for \textbf{communities} \textbf{facing} \textbf{environmental} harm, and her work illustrates how litigation and legal advocacy can confront \textbf{corporations} and \textbf{institutions} that \textbf{profit} from pollution and \textbf{environmental} \textbf{destruction}.
Liu Junyan focuses on \textbf{climate} \textbf{risks} and \textbf{inequality}, emphasizing that \textbf{climate} change does not affect everyone equally.
By drawing attention to who is most exposed to \textbf{climate} \textbf{impacts} and who has the least \textbf{resources} to respond, Liu ’s work underlines that \textbf{climate} \textbf{policy} must tackle structural \textbf{inequality} if it is to be truly effective and just.
Pua Lay Peng is fighting plastic waste, confronting one of the most visible and pervasive forms of pollution.
Her efforts show how tackling plastic—from production to disposal—is not only about cleaning up trash, but also about protecting \textbf{communities}, \textbf{ecosystems}, and \textbf{future} \textbf{generations} from toxic and long-lasting \textbf{impacts}.
Lawyer Sylvia Wu works on \textbf{climate} litigation and fishery \textbf{rights}, connecting the dots between \textbf{climate} change, ocean health, and the \textbf{livelihoods} of coastal and fishing \textbf{communities}.
Her work reinforces that \textbf{climate} \textbf{justice} includes defending the \textbf{rights} of those who depend directly on the sea, and that legal strategies can be a powerful way to protect both \textbf{people} and marine \textbf{ecosystems}.
Community organizer Veronica Cabe demonstrates how local organizing can transform \textbf{climate} and \textbf{environmental} \textbf{struggles}.
By bringing \textbf{people} together, building collective \textbf{power}, and amplifying \textbf{community} \textbf{voices}, she shows that \textbf{climate} action \textbf{rooted} in local \textbf{realities} is both more democratic and more resilient.
Yaewon Hwang ’s work on trans \textbf{justice} makes clear that \textbf{climate} \textbf{justice} is inseparable from \textbf{gender} \textbf{justice}.
Trans and \textbf{gender} non-conforming \textbf{people} often \textbf{face} heightened discrimination and \textbf{risk} in times of \textbf{crisis}, from \textbf{displacement} to access to services. \textbf{Centering} trans \textbf{justice} in \textbf{climate} work \textbf{challenges} exclusionary approaches and insists that safety, \textbf{dignity}, and \textbf{rights} for all \textbf{genders} are non-negotiable in a just transition.
Ai Ji focuses on \textbf{biodiversity} threats, reminding us that the \textbf{climate} \textbf{crisis} and the \textbf{biodiversity} \textbf{crisis} are deeply intertwined.
The \textbf{loss} of species and \textbf{ecosystems} not only accelerates \textbf{climate} \textbf{breakdown}, it also undermines the \textbf{resilience} of \textbf{communities} who depend on healthy \textbf{forests}, oceans, and landscapes.
Across these diverse struggles—Indigenous \textbf{forest} \textbf{protection}, anti-coal \textbf{resistance}, legal fights against polluters, \textbf{challenges} to \textbf{climate} \textbf{inequality}, campaigns against plastic, \textbf{defense} of fisheries, \textbf{grassroots} organizing, trans \textbf{justice}, and \textbf{biodiversity} protection—runs a common thread : the \textbf{leadership} of \textbf{women} and \textbf{gender} non-conforming \textbf{people}, and the \textbf{demand} for \textbf{climate} action that is inclusive, rights-based, and intersectional.
To move forward, \textbf{environmental} movements must confront eco-ableism, dismantle \textbf{barriers} that exclude disabled \textbf{people}, and actively create \textbf{space} for the \textbf{leadership} of those most affected by the \textbf{climate} \textbf{crisis}.
Listening to and \textbf{supporting} activists like Archana Soreng, Eunbin Kang, Hasminah D.
Paudac-Tawano, Liu Junyan, Pua Lay Peng, Sylvia Wu, Veronica Cabe, Yaewon Hwang, and Ai Ji is not optional—it is fundamental to building the just, livable \textbf{future} that \textbf{climate} action is supposed to deliver.

% matched lemmas: barrier, biodiversity, breakdown, center, central, challenge, climate, coal, community, corporation, crisis, defense, demand, destruction, dignity, displacement, ecosystem, environmental, face, force, forest, frontline, fuel, future, gender, generation, grassroots, guardian, impact, inequality, institution, justice, knowledge, leadership, livelihood, loss, mean, narrative, people, policy, power, priority, profit, protection, push, reality, resilience, resistance, resource, right, risk, root, space, struggle, support, system, tool, voice, woman, world, youth
\end{textsample}
